\documentclass[11pt]{amsart}
\usepackage[T1]{fontenc}
\usepackage{lmodern}
\usepackage[letterpaper,margin=1in]{geometry}
\usepackage{microtype}
\usepackage{mathtools,amssymb}
\usepackage{xcolor}
\definecolor{linkblue}{rgb}{0,0,0.6}
\usepackage[colorlinks=true,linkcolor=linkblue,citecolor=linkblue,urlcolor=linkblue]{hyperref}
\hypersetup{pdftitle={Infinitely Many Decreases in Distinct-Power Thresholds},
pdfauthor={Will Howe; Grant Huffman; Jensen Kohlmeyer; Liam Kruer}}
\numberwithin{equation}{section}
\newtheorem{theorem}{Theorem}[section]
\newtheorem{proposition}[theorem]{Proposition}
\newtheorem{lemma}[theorem]{Lemma}
\newtheorem{corollary}[theorem]{Corollary}
\theoremstyle{remark}
\newtheorem{remark}[theorem]{Remark}
\newcommand{\N}{\mathbb N}
\newcommand{\Z}{\mathbb Z}
\newcommand{\Q}{\mathbb Q}
\newcommand{\R}{\mathbb R}
\newcommand{\C}{\mathbb C}
\newcommand{\e}{\mathrm e}
\newcommand{\ii}{\mathrm i}
\newcommand{\E}{\mathbb E}
\newcommand{\Prob}{\mathbb P}
\newcommand{\normone}[1]{\lVert #1\rVert_1}
\newcommand{\dist}[1]{\lVert #1\rVert_{\R/\Z}}
\DeclareMathOperator{\supp}{supp}
\DeclareMathOperator{\Var}{Var}
\setlength{\emergencystretch}{2em}
\allowdisplaybreaks[1]
\title[Distinct-power thresholds]{Infinitely many decreases in the thresholds\\
for sums of distinct powers}
\author[W. Howe]{Will Howe}
\author[G. Huffman]{Grant Huffman}
\author[J. Kohlmeyer]{Jensen Kohlmeyer}
\author[L. Kruer]{Liam Kruer}
\date{November 8, 2026}
\subjclass[2020]{11P05, 05A17}
\keywords{distinct powers, subset sums, complete sequences, cyclotomic recurrences}
\thanks{Email addresses:
\href{mailto:howew@purdue.edu}{\texttt{howew@purdue.edu}} (Will Howe);
\href{mailto:huffma97@purdue.edu}{\texttt{huffma97@purdue.edu}} (Grant Huffman);
\href{mailto:jenwin@purdue.edu}{\texttt{jenwin@purdue.edu}} (Jensen Kohlmeyer);
\href{mailto:kruerl@purdue.edu}{\texttt{kruerl@purdue.edu}} (Liam Kruer).}
\begin{document}
\begin{abstract}
Let \(T_k\) be the least integer such that every integer at least \(T_k\)
is a sum of distinct positive \(k\)th powers. We prove that
\(T_k>T_{k+1}\) for infinitely many positive integers \(k\).
The principal upper estimate holds on a set of degrees meeting every
sufficiently large dyadic interval:
\[
 \log T_k\le(1+o(1))\,k
 \exp\!\left(\frac25\frac{\log k}{\log\log k}\right).
\]
We obtain it from flat integral recurrences, cyclotomic coefficient
estimates, an exact reduction of polynomial degree, and a subset-sum
construction. A shifted-prime divisor theorem of Fan and Pollack supplies
larger lower bounds along an unbounded sequence of degrees. The two
estimates contradict eventual monotonicity of the thresholds.
\end{abstract}
\maketitle

\section{The problem and the two estimates}
For \(A\subseteq\Z_{\ge0}\), write
\[
 \Sigma(A)=\left\{\sum_{a\in B}a:B\subseteq A,\ B\text{ finite}\right\}.
\]
The empty sum is allowed. Let \(\mathcal P_k=\{n^k:n\in\N,\ n\ge1\}\), and put
\[
 T_k=\min\{T\in\Z_{\ge0}:[T,\infty)\cap\Z\subseteq\Sigma(\mathcal P_k)\}.
\]
The finiteness of \(T_k\) for every positive \(k\) is classical; see
Kim~\cite{Kim}. We have \(T_1=0\). If \(k\ge2\), the integer two is
not representable, so \(T_k\ge3\). For \(k\ge2\), using the largest
exceptional integer instead subtracts one from each threshold, and
therefore preserves adjacent comparisons. The question considered here
is Erd\H{o}s problem~345; see Erd\H{o}s and Graham~\cite[p.~55]{EG}.

\begin{theorem}\label{thm:main}
There are infinitely many positive integers \(k\) for which
\(T_k>T_{k+1}\).
\end{theorem}

All logarithms are natural, and
\(\mathcal L(x)=\log x/\log\log x\) for \(x>\e\).
An unspecified positive absolute constant may be enlarged when it
reappears; dependence on a fixed parameter is indicated by a subscript.
All asymptotic estimates below are uniform in the degrees and
frequencies to which they are applied.

\begin{proposition}\label{prop:upper}
There is a set \(\mathcal G\subseteq\N\) meeting every sufficiently large
interval \([2^m,2^{m+1}]\) such that
\begin{equation}\label{eq:upper}
 \log T_k\le(1+o(1))\,k\exp\!\left(\frac25\mathcal L(k)\right)
 \qquad(k\in\mathcal G).
\end{equation}
\end{proposition}
This is proved in Sections~\ref{sec:flat}--\ref{sec:completion}.
For comparison, define
\[
 \omega^*(n)=\#\{p:p\text{ prime},\ p-1\mid n\},\qquad
 D(n)=\prod_{\substack{p\text{ prime}\\p-1\mid n}}p.
\]
Fan and Pollack~\cite[Theorem~1]{FP} prove unconditionally that
\begin{equation}\label{eq:FP}
 \omega^*(n)>
 \exp\!\left((0.6736\log2)\mathcal L(n)\right)
\end{equation}
for infinitely many \(n\). Only this infinitely-often conclusion is used.

\begin{lemma}\label{lem:lower}
There are arbitrarily large \(j\) such that
\begin{equation}\label{eq:lower}
 \log T_j\ge j\exp\!\left(\frac9{20}\mathcal L(j)\right).
\end{equation}
\end{lemma}
\begin{proof}
Choose \(n\) satisfying \eqref{eq:FP}, and put \(j=n+1\) and \(Q=D(n)\).
For every integer \(b\), \(b^j\equiv b\pmod Q\). Modulo a prime factor
\(p\) of \(Q\), this follows from \(p-1\mid n\) if \(p\nmid b\);
if \(p\mid b\), both sides vanish. The factors of \(Q\) are distinct.

Set \(H=\lfloor\sqrt Q\rfloor\) and \(S=H(H+1)/2<Q\).
Subset sums of powers with bases at most \(H\) have residues in
\(\{0,\ldots,S\}\). Choose \(M\equiv S+1\pmod Q\) with
\[
 (H+1)^j-Q\le M<(H+1)^j.
\]
For large \(Q\) and \(j\ge3\), this is positive. It cannot be represented:
a base larger than \(H\) contributes too much, while the smaller bases
cannot produce its residue. Hence
\[
 \log T_j\ge\log M\ge \frac j2\log Q-O(j).
\]
Since \(\log Q\ge(\log2)\omega^*(n)\),
\(0.6736\log2>9/20\), and
\(\mathcal L(n+1)-\mathcal L(n)\to0\), the fixed factors and the error
are absorbed by the strict difference of constants. This proves
\eqref{eq:lower}.
\end{proof}

\section{Flat products and integral recurrences}\label{sec:flat}
For \(P(z)=\sum_h w_hz^h\), put \(\normone P=\sum_h|w_h|\).
Let \(E\) denote the shift operator on sequences:
\((Ex)_n=x_{n+1}\), and write \(\Delta=E-1\).
A rational polynomial is called integer-valued if its values at all
integers are integers.

We use the following consequence of Belov and
Konyagin~\cite[Corollary~0.2]{BK}. For every positive integer \(J\), there
are nonincreasing nonnegative integers \(a_1,\ldots,a_J\), of sum \(J\),
and a real \(a_0\), such that
\begin{equation}\label{eq:cosine}
 a_0+\sum_{h=1}^J a_h\cos(h\theta)\ge0\quad(\theta\in\R),
 \qquad a_0\le2^{32}(1+\log J)^5.
\end{equation}

\begin{lemma}\label{lem:flat}
For every \(J\ge1\), one can choose
\[
 P_J(z)=\prod_{h=1}^J(1-z^h)^{a_h}=(1-z)^JQ_J(z)\in\Z[z]
\]
so that
\begin{align}
 \deg P_J&\le J^2,&
 \log\normone{P_J}&\le C(1+\log J)^6,\label{eq:flatnorm}\\
 Q_J&\in\Z_{\ge0}[z],&
 M_J:=Q_J(1)&=\prod_h h^{a_h}\mid J!,\label{eq:moment}\\
 \sum_{\substack{h\le J\\v\mid h}}a_h&\le J/v&& (v\ge1).
 \label{eq:slots}
\end{align}
The constant and leading coefficients of \(P_J\) have absolute value one.
\end{lemma}
\begin{proof}
Take the coefficients from \eqref{eq:cosine}. For \(0<\rho<1\),
\[
 \log|P_J(\rho\e^{\ii\theta})|
 =-\sum_{m\ge1}\frac{\rho^m}{m}\sum_h a_h\cos(hm\theta)
 \le a_0\log\frac1{1-\rho}.
\]
Also
\[
 |1-\e^{\ii t}|\le\frac2{1+\rho}|1-\rho\e^{\ii t}|.
\]
Taking \(\rho=1-J^{-1}\), for \(J\ge2\), gives
\(\log\|P_J\|_\infty\le C(1+\log J)^6\).
By Parseval,
\(\normone{P_J}\le\sqrt{\deg P_J+1}\|P_J\|_\infty\), proving
\eqref{eq:flatnorm}; \(J=1\) is immediate.

The product
\[
 Q_J(z)=\prod_h(1+z+\cdots+z^{h-1})^{a_h}
\]
has nonnegative integral coefficients.
Let \(\lambda_i=\#\{h:a_h\ge i\}\), \(1\le i\le a_1\).
Then \(\sum_i\lambda_i=J\) and
\(M_J=\prod_i\lambda_i!\), which divides \(J!\) by the multinomial
formula. Finally, the block of \(v\) indices ending at \(iv\) has
coefficient sum at least \(va_{iv}\). Adding the disjoint blocks proves
\eqref{eq:slots}. The remaining assertions follow from the product.
\end{proof}

\begin{lemma}[Rounding to an exact recurrence]\label{lem:rounding}
Let \(k,X\) be positive integers and put \(K=1+\log k\).
Suppose \(1\le d\le k\), \(F\in\R[x]\) has degree at most \(d\), and
\[
 \sum_{0\le n<X}\sin^2(\pi F(n))\le k^2.
\]
If \(\log X\) grows faster than every fixed power of \(K\), there is an
integer sequence \(m_n\), defined on all of \(\Z\), with
\(P_{d+1}(E)m=0\), and an interval \(I\) of \(L\) consecutive integers
such that
\begin{equation}\label{eq:roundenergy}
 \sum_{n\in I}|F(n)-m_n|^2\le k^2/4,\qquad
 L\ge X\exp(-CK^6).
\end{equation}
\end{lemma}
\begin{proof}
On the prefix choose \(m_n\) nearest to \(F(n)\), and put
\(\epsilon_n=F(n)-m_n\).
Since \(\sin(\pi t)\ge2t\) on \(0\le t\le1/2\),
\(\sum\epsilon_n^2\le k^2/4\).
The root of \(P_{d+1}\) at one has order \(d+1\), so
\(P_{d+1}(E)F=0\).
The residuals \(P_{d+1}(E)m\) are integers, and the finite convolution
inequality gives
\[
 \sum_n|P_{d+1}(E)m_n|^2
 \le\normone{P_{d+1}}^2\sum_n\epsilon_n^2
 \le\normone{P_{d+1}}^2 k^2/4.
\]
Write \(D_0=\deg P_{d+1}\). The valid residual positions are
\(0,\ldots,X-D_0-1\). If \(B_0\) is the number of nonzero residuals,
then \(B_0\le \normone{P_{d+1}}^2k^2/4\).
The zero residuals split into at most \(B_0+1\) runs, so one run has
length at least
\[
 \frac{X-D_0-B_0}{B_0+1}.
\]
Here \(D_0\le(d+1)^2\) and \(B_0+1\le\exp(CK^6)\).
The hypothesis on \(X\) therefore gives a run of \(r\) zero residuals,
beginning at an integer \(s\), with \(r\ge X\exp(-CK^6)\), after
enlarging \(C\).

Use the \(D_0\) rounded values at \(s,\ldots,s+D_0-1\) to define an
exact recurrence sequence. The leading coefficient of \(P_{d+1}\)
is a unit, so the \(r\) zero residual equations force agreement with
the rounded values throughout \(s,\ldots,s+r+D_0-1\).
Its constant coefficient is also a unit, so the extension to all of
\(\Z\) is integral in both directions.
This whole interval lies in the original prefix. Restricting the
error estimate to it proves the assertion.
\end{proof}

\section{Cyclotomic heights on a long interval}\label{sec:heights}
Write
\[
 P(z)=\prod_{i=1}^J(1-z^{h_i}),\qquad
 1\le h_i\le H,\qquad D=\deg P.
\]
Every sequence annihilated by \(P(E)\) is a sum
\(\sum_\zeta p_\zeta(n)\zeta^{-n}\), where the roots of unity have
orders at most \(H\), and the degree of each polynomial is less than
the multiplicity of its root.

\begin{lemma}[Long-interval separation]\label{lem:gram}
For such a sum \(x_n\), let \(a_\zeta\) be the leading coefficient of
a nonzero component of degree \(b\). On an interval of \(L\) consecutive
integers,
\begin{equation}\label{eq:gram}
 \sum_{n\in I}|x_n|^2
 \ge\frac{L^{2b+1}|a_\zeta|^2}
 {2(2b+1)\binom{2b}{b}^2},
\end{equation}
provided \(L\ge CD(J+1)^3H^2\).
\end{lemma}
\begin{proof}
Translate the interval to start at zero.
Let \(\ell_j\) be the shifted Legendre polynomials orthonormal on
\([0,1]\), with leading coefficient
\(\sqrt{2j+1}\binom{2j}{j}\).
Writing \(\ell_j(x)=\sqrt{2j+1}\,\mathsf P_j(2x-1)\) in terms of the usual
Legendre polynomial, one has
\[
 \|\ell_j\|_\infty\le\sqrt{2j+1},\qquad
 \|\ell_j'\|_\infty\le j(j+1)\sqrt{2j+1}.
\]
For example, the derivative bound follows by summing
\(\mathsf P_j'=\sum_{0\le r<j,\ j-r\ {\rm odd}}(2r+1)\mathsf P_r\) and using
\(|\mathsf P_r|\le1\) on \([-1,1]\).
Products of two basis polynomials of degrees less than \(J\)
therefore have supremum and total variation \(O((J+1)^3)\).

Use the basis \(\ell_j(n/L)\zeta^{-n}\), with the appropriate degrees
at each root. Its dimension is at most \(D\).
At a common frequency, the normalized discrete inner product differs
from the continuous one by \(O((J+1)^3/L)\).
Distinct root frequencies are separated modulo one by at least \(H^{-2}\).
Summation by parts and the geometric-series bound therefore give
off-diagonal entries \(O((J+1)^3H^2/L)\).
The row-sum estimate shows that the normalized Gram matrix has all
eigenvalues at least \(1/2\) under the stated length condition.
The coefficient of the top basis vector in the chosen component is
\[
 \frac{a_\zeta L^b}{\sqrt{2b+1}\binom{2b}{b}}.
\]
Its contribution to the Gram lower bound gives \eqref{eq:gram}.
Translation multiplies the top coefficient by a number of absolute
value one and does not change this estimate.
\end{proof}

\begin{lemma}[Integer recurrence height]\label{lem:integerheight}
If a nonzero integral sequence \(x\) is annihilated by \(P(E)\), and
\(b\) is its largest component degree, then
\begin{equation}\label{eq:integerheight}
 (2b+1)\log L\le
 \log\!\left(\sum_{n\in I}|x_n|^2\right)
 +CJ\log(2H(J+1))
\end{equation}
whenever \(L\) satisfies Lemma~\ref{lem:gram}.
\end{lemma}
\begin{proof}
Translate the interval to start at zero.
The recurrence first gives a generating-function denominator equal
to the reciprocal of \(P\). Since
\(z^D P(1/z)=(-1)^J P(z)\), a change of sign in the numerator writes
the generating function as \(A/P\), with \(A\in\Z[z]\) of degree less
than \(D\).
Cancel the monic common cyclotomic factors. Gauss's lemma gives
\(A_0/C\) with relatively prime integral polynomials. Choose a primitive root \(\zeta\) at which
\(C\) has multiplicity \(b+1\).
The absolute value of the top component coefficient is
\[
 \frac{|A_0(\zeta)|}{|C_{\rm reg}(\zeta)|\,b!},
 \qquad C_{\rm reg}(z)=C(z)/(z-\zeta)^{b+1}.
\]
The numerator is a nonzero algebraic integer, and its geometric mean
over primitive conjugates is at least one.

If \(P=GC\), then \(P_{\rm reg}=G_{\rm reg}C_{\rm reg}\).
Regularized values of integral polynomials are algebraic integers:
successive division by the monic factor \(z-\zeta\) preserves algebraic
integrality. Thus the geometric mean of \(|G_{\rm reg}(\zeta)|\)
is at least one. Each vanishing binomial in \(P\) contributes a derivative
of absolute value at most \(H\), and each nonvanishing one contributes
at most two. The geometric mean of \(|C_{\rm reg}(\zeta)|\) is consequently
at most \((2H)^J\). Some conjugate top coefficient has absolute value
at least \(((2H)^J J!)^{-1}\).
Insert it into \eqref{eq:gram}, use \(\binom{2b}{b}\le4^b\), and
take logarithms.
\end{proof}

\begin{lemma}[Refined cyclotomic height]\label{lem:refined}
Let \(K=1+\log k\), \(Z=\log X\), with \(Z\) growing faster than every
fixed power of \(K\). Suppose \(2\le q\mid k!(k+1)!\), \(\gcd(u,q)=1\),
and \(F\in\Q[x]\) has degree at most \(d\), \(1\le d\le k\), with
\[
 qF(n)\in\Z,\qquad qF(n)\equiv un^k\pmod q\quad(n\in\Z).
\]
If its sine energy on \(0\le n<X\) is at most \(k^2\), there is a
nonzero rational sequence \(\epsilon\), annihilated by \(P_{d+1}(E)\),
such that \(q\epsilon\) is integral and congruent to \(un^k\pmod q\)
globally. If \(h_t\) denotes its multiplicity at primitive \(t\)th
roots, then
\begin{align}
 h_t&\le\frac{CdK}{tZ}\quad(t\ge2),\label{eq:nontrivialheight}\\
 h_1&\le\frac{C(\log q+dK)}Z.\label{eq:oneheight}
\end{align}
In particular a multiplicity vanishes if its upper bound is less than one.
\end{lemma}
\begin{proof}
Apply Lemma~\ref{lem:rounding} and put \(\epsilon=F-m\).
This gives all the integrality, congruence, and energy assertions,
with energy at most \(k^2/4\) on an interval of length
\(L\ge X\exp(-CK^6)\). The sequence is nonzero by its congruence at
\(n=1\). Its least annihilator is an integral product of cyclotomic
factors dividing \(P_{d+1}\), by cancellation in the generating
function of \(q\epsilon\).

Put \(J=d+1\). At a nontrivial root the components of \(\epsilon\)
are the negatives of those of the integer sequence \(m\), since \(F\)
contributes only at one. Fix primitive order \(t\ge2\) with
\(h_t=b+1>0\). We refine the denominator estimate in the preceding proof.
There are at most \(J/t\) vanishing factors, by \eqref{eq:slots}, and
their derivatives have absolute value at most \(J\).

For a nonvanishing factor \(1-\zeta^h\), put \(l=t/\gcd(t,h)>1\).
Its geometric mean over primitive \(t\)th roots is
\(\Phi_l(1)^{1/\varphi(l)}\).
This is one unless \(l=p^a\) is a prime power, when its logarithm is
\(\log p/\varphi(p^a)\). Such factors have \(t/p^a\mid h\);
by \eqref{eq:slots} their number with multiplicity is at most \(Jp^a/t\).
The total logarithmic contribution of nonvanishing factors is thus
at most
\[
 \sum_{p^a\mid t}\frac{Jp^a}t
       \frac{\log p}{\varphi(p^a)}
 \le\frac{2J}t\sum_{p^a\mid t}\log p
 =\frac{2J}t\log t.
\]
Cancellation decreases the geometric-mean denominator bound, as above,
and \(b!\le J^{J/t}\). Some conjugate top coefficient therefore satisfies
\[
 |a_\zeta|\ge\exp\!\left(-C(J/t)\log(2J)\right).
\]
The interval-length condition in Lemma~\ref{lem:gram} is polynomial in
\(k\), and hence holds here. It follows that
\[
 (2h_t-1)\log L\le C(J/t)K+C\log k.
\]
If \(h_t>0\), then \(t\le J\), so \(J/t\ge1\).
Now \(2h_t-1\ge h_t\), \(J\le2d\), and \(\log L\sim Z\).
This proves \eqref{eq:nontrivialheight} without an additive constant.

Apply Lemma~\ref{lem:integerheight} to \(q\epsilon\) for the remaining
bound. Its energy is at most \(q^2k^2/4\), and its largest component
degree is at least \(h_1-1\). This gives \eqref{eq:oneheight};
in particular the term involving \(\log q\) is not discarded.
\end{proof}

\section{Exact arithmetic of polynomial reductions}\label{sec:arithmetic}

For \(0\le d<k\), define
\[
 A_j(k)=\Delta^j(x^k)\big|_{x=0}=j!\,S(k,j),\qquad
 G_d(k)=\gcd\{A_j(k):d<j\le k\},
\]
where \(S(k,j)\) is a Stirling number of the second kind.
In particular \(G_0(k)=1\), since \(A_1(k)=1\).
The Newton expansion shows that
\begin{equation}\label{eq:newtoncriterion}
 q\mid G_d(k)
 \quad\Longleftrightarrow\quad
 \begin{gathered}
 \text{there exists an integer-valued }P\in\Q[x],\ \deg P\le d,\\
 n^k\equiv P(n)\pmod q\quad\text{for every }n\in\Z.
 \end{gathered}
\end{equation}
Indeed, the forward differences above degree \(d\) must vanish modulo
\(q\); conversely, the truncation
\(\sum_{j=0}^d A_j(k)\binom{x}{j}\) supplies \(P\).
This formulation applies on all integers, including negative ones.

\begin{lemma}[A product bound for the reduction modulus]\label{lem:product}
For \(1\le d<k\), put
\[
 F_{k,j}=\gcd_{a\in\Z_{\ge0}}(a^k-a^j)\quad(1\le j\le d).
\]
Then
\begin{equation}\label{eq:product}
 G_d(k)\mid\prod_{j=1}^d F_{k,j},\qquad
 F_{k,j}\mid 2(k-j)D(k-j).
\end{equation}
\end{lemma}
\begin{proof}
Start with \(q_d=G_d(k)\), and recursively put
\(q_{j-1}=\gcd(q_j,A_j(k))\).
Then \(q_j\mid G_j(k)\), and the canonical truncation
\[
 P_j(x)=\sum_{\ell=0}^j A_\ell(k)\binom{x}{\ell}
\]
is congruent to \(x^k\pmod {q_j}\) on the integers.
For every nonnegative integer \(a\), the polynomial
\(P_j(ax)-a^kP_j(x)\) is integer-valued and divisible by \(q_j\) at
every integer. Its \(j\)th forward difference is
\(A_j(k)(a^j-a^k)\). Thus
\[
 q_j/q_{j-1}\mid F_{k,j}.
\]
Since \(q_0=1\), multiplication gives the first assertion.

Write \(h=k-j>0\). A prime \(p\) divides \(F_{k,j}\) exactly when
\(p-1\mid h\). Necessity follows by testing units modulo \(p\);
sufficiency also includes nonunits, because \(j\ge1\).
Testing \(a=p\) bounds \(v_p(F_{k,j})\) by \(j\).
For odd \(p\), testing \(a=1+p\) bounds it by \(1+v_p(h)\).
For \(p=2\), testing \(a=3\) bounds it by \(2+v_2(h)\).
These lifting identities show that \(F_{k,j}\mid2hD(h)\).
Taking logarithms gives
\begin{equation}\label{eq:logG}
 \log G_d(k)\le d\log(2k)+\sum_{j=1}^d\log D(k-j).
\end{equation}
\end{proof}

We next control a prime-power modulus from the cyclotomic multiplicities
of an annihilator. Here \(v_p\) denotes the \(p\)-adic valuation.

\begin{lemma}[Dilation of an integral recurrence]\label{lem:dilation}
Let \(p\) be prime, \(e\ge1\), and
\[
 B(z)=(z-1)^{m_0}\prod_{a=1}^A\Phi_{p^a}(z)^{m_a},
 \qquad 0\le m_a\le k\quad(0\le a\le A).
\]
If \(B(E)n^k\equiv0\pmod {p^e}\) for all \(n\in\Z\), then
\begin{equation}\label{eq:dilation}
 e\le C\left\{\frac{m_0}{p-1}
       +\sum_{a=1}^A(a+1)m_a
       +(A+1)(1+\log_p(k+1))\right\}.
\end{equation}
\end{lemma}
\begin{proof}
Put \(D_B=\deg B\). The rational \(B\)-recurrence sequences on
\(\Z\) have an integral lattice identified with \(\Z^{D_B}\) by their
first \(D_B\) entries, since the leading and constant coefficients
of \(B\) are units.
For a positive integer \(g\) prime to \(p\), the dilation
\(S_g(x)_n=x_{gn}\) preserves this space and its integral lattice:
on the complex basis \(n^j\zeta^n\), it permutes primitive roots of a
fixed \(p\)-power order and multiplies the degree-\(j\) vector by \(g^j\).

Extend \(0^k,1^k,\ldots,(D_B-1)^k\) by the recurrence to an integral
sequence \(y\). The assumed congruence makes \(y_n\equiv n^k\pmod {p^e}\)
for all integers \(n\). The cases \(D_B=0,1\) are impossible, because
the values at zero and one are respectively zero and one.
Thus the initial vector contains a unit, and
\[
 S_g y\equiv g^k y\pmod {p^e}.
\]
If \(\chi_g\) is the characteristic polynomial of this integral
matrix, Cayley--Hamilton gives
\begin{equation}\label{eq:characteristic}
 p^e\mid\chi_g(g^k).
\end{equation}

For odd \(p\), choose \(g\) primitive modulo \(p^2\), hence modulo
every \(p^a\). The primitive \(p^a\)th roots form one cycle of length
\(\varphi(p^a)\). Their degree-\(j\) characteristic factor is
\(T^{\varphi(p^a)}-g^{j\varphi(p^a)}\), for \(0\le j<m_a\).
At \(T=g^k\), its valuation is
\[
 a+v_p(k-j).
\]
The root-one factor \(T-g^j\) has valuation zero unless \(p-1\mid k-j\);
in that case its valuation is \(1+v_p(k-j)\).
None of these factors vanishes, since \(j<m_a\le k\).

For any interval of \(m\le k\) consecutive positive integers ending
at \(k\), counting multiples of \(p^\nu\) gives
\begin{equation}\label{eq:valuationinterval}
 \sum_{j=0}^{m-1}v_p(k-j)
 \le \frac m{p-1}+\log_p(k+1).
\end{equation}
For the root-one contribution, restrict to the single residue class
\(k-j\equiv0\pmod {p-1}\), and divide its terms by \(p-1\).
The resulting consecutive integers number at most \(m_0/(p-1)+1\).
Applying \eqref{eq:valuationinterval} gives a contribution at most
\(C m_0/(p-1)+C(1+\log_p(k+1))\).
The other cycles contribute at most
\(\sum_a(a+1)m_a+A\log_p(k+1)\).
Equation~\eqref{eq:characteristic} proves the claim for odd \(p\).

For \(p=2\), take \(g=5\). There are at most two cycles at each
primitive \(2^a\)th order, of length \(2^{a-2}\) for \(a\ge3\);
the orders two and four are handled directly.
The valuation of each factor is at most
\(C(a+1+v_2(k-j))\).
The root-one contribution is \(2+v_2(k-j)\).
Summing with \eqref{eq:valuationinterval} proves the same bound.
\end{proof}

\begin{lemma}[One contraction step]\label{lem:contraction}
Under the hypotheses of Lemma~\ref{lem:refined}, if the sine energy is
at most \(k^2\), there is an integer \(M\ge1\) satisfying
\begin{equation}\label{eq:contraction}
 M\le \frac{CK^2(\log q+dK)}Z+CK^4
\end{equation}
such that, whenever \(M<k\), one has \(q\mid G_M(k)\).
All constants are uniform for \(1\le d\le k\).
\end{lemma}
\begin{proof}
Choose
\[
 T=C_0\frac{\log q+dK}Z
\]
large enough that Lemma~\ref{lem:refined} gives \(h_1\le T\) and
\(h_{p^a}\le T/p^a\). These multiplicities are less than \(k\) for
large \(k\), since \(\log q=O(kK)\) and \(Z\) dominates powers of \(K\).

Fix \(p^e\Vert q\). Factor the integral annihilator of \(\epsilon\)
as \(BC\), where \(B\) contains precisely its root-one and
\(p\)-power cyclotomic factors. Then \(C(1)\) is a unit modulo \(p\):
\(\Phi_t(1)\) is one unless \(t\) is a prime power, and the prime-power
factors in \(C\) correspond to primes other than \(p\).
Put \(c_0=C(1)\) and write \(C(1+\Delta)=c_0+\Delta R(\Delta)\).
On polynomial sequences of degree at most \(k\) modulo \(p^e\),
\(\Delta^{k+1}=0\). An explicit inverse is therefore
\[
 c_0^{-1}\sum_{j=0}^{k}
       \bigl(-c_0^{-1}\Delta R(\Delta)\bigr)^j.
\]
Because \(B(E)C(E)(q\epsilon)=0\) and \(q\epsilon\equiv un^k\pmod {p^e}\),
the unit \(u\) may be cancelled.
Apply the displayed inverse, which commutes with \(B(E)\), to obtain
\[
 B(E)n^k\equiv0\pmod {p^e}.
\]
Apply Lemma~\ref{lem:dilation}. The largest \(a\) is at most
\(\log_p(d+1)\), and
\(\sum_{a\ge1}(a+1)p^{-a}=O((p-1)^{-1})\). Hence
\begin{equation}\label{eq:primeexponent}
 e\le CT/(p-1)+CK^2.
\end{equation}
If \(p>T\), all nontrivial \(p\)-power roots are absent.
Then \(B=(z-1)^{h_1}\); necessarily \(h_1\ge1\), and
\[
 p^e\mid G_{h_1-1}(k).
\]

Set \(M=\lceil C_1(T+K^2)K^2\rceil\), with \(C_1\) fixed and sufficiently
large. For \(p\le T\), the elementary bound
\(v_p(M!)\ge M/(2p)\), valid here since \(M\ge2p\), dominates
\eqref{eq:primeexponent}. Thus the part of \(q\) supported on primes
at most \(T\) divides \(M!\).
Now assume \(M<k\), as in the conclusion.
The remaining prime powers divide \(G_M(k)\), because \(M\ge h_1\)
and \(G_{h_1-1}(k)\mid G_M(k)\).
Each \(A_j(k)=j!S(k,j)\), \(j>M\), is divisible by \(M!\).
Thus the small-prime part also divides \(G_M(k)\).
The two parts are coprime, so \(q\mid G_M(k)\).
The chosen value of \(M\) satisfies \eqref{eq:contraction}.
\end{proof}

\section{Good degrees and the exact-center estimate}\label{sec:centers}

Fix throughout
\begin{equation}\label{eq:parameters}
 c=\frac25,\qquad a=\frac1{10},\qquad \beta=\frac1{100}.
\end{equation}

\begin{lemma}[Simultaneous degree selection]\label{lem:selection}
For every sufficiently large dyadic \(Y\), all but \(o(Y)\) integers
\(k\in[Y,2Y]\) satisfy
\begin{equation}\label{eq:good}
 \log D(k-j)\le \exp(a\mathcal L(Y))
 \quad\text{for }1\le j\le\lceil(2Y)^\beta\rceil.
\end{equation}
\end{lemma}
\begin{proof}
Let \(\tau\) be the divisor function and \(d_z\) the ordered
\(z\)-fold divisor function. For every integer \(m\ge1\),
\[
 \tau(n)^m\le d_{2^m}(n).
\]
For a prime power exponent \(b\), this follows by induction from
\[
 \frac{\binom{b+2z-1}{b}}{\binom{b+z-1}{b}}
 =\prod_{i=0}^{b-1}\frac{2z+i}{z+i}\ge b+1\qquad(z\ge1);
\]
each factor increases with \(z\), and the product at \(z=1\) is \(b+1\).
Multiplicativity completes the proof.
Counting ordered factorizations and summing the first \(z-1\) factors
harmonically gives, with no constant depending on \(z\),
\begin{equation}\label{eq:momentbound}
 \sum_{n\le U}d_z(n)\le U(1+\log U)^{z-1}.
\end{equation}
Since a prime predecessor divisor of \(n\) is at most \(n+1\),
\(\log D(n)\le\tau(n)\log(n+1)\).

Take
\[
 m=\left\lfloor\frac{\log\log Y-3\log\log\log Y}{\log2}\right\rfloor,
 \qquad
 B_Y=\frac{\exp(a\mathcal L(Y))}{\log(2Y+1)}.
\]
By \eqref{eq:momentbound} and Markov's inequality, the number of
\(n\le2Y\) with \(\log D(n)>\exp(a\mathcal L(Y))\) is at most
\[
 2Y(1+\log(2Y))^{2^m-1}B_Y^{-m}
   =Y^{1-a/\log2+o(1)}.
\]
Indeed \(2^m\le\log Y/(\log\log Y)^3\), so the positive logarithmic
error is \(o(\log Y)\), while
\(m\log B_Y=(a/\log2+o(1))\log Y\).
The union over at most \(O(Y^\beta)\) shifts has size
\(Y^{1-a/\log2+\beta+o(1)}=o(Y)\), because
\(\beta<a/\log2\). All shifted arguments are positive for large \(Y\).
\end{proof}

Let \(\mathcal G\) be the union of the sets selected in
Lemma~\ref{lem:selection}, and for \(k\in\mathcal G\) put
\begin{equation}\label{eq:NX}
 K=1+\log k,\qquad
 N=\left\lceil\exp\bigl(\exp(c\mathcal L(k))\bigr)\right\rceil,\qquad
 X=\lfloor N\exp(-B K^6)\rfloor ,
\end{equation}
where \(B>0\) is any fixed constant. Then
\begin{equation}\label{eq:scales}
 Z:=\log X\sim\log N\sim\exp(c\mathcal L(k)),\qquad
 K^A=o(Z)\ (A\text{ fixed}),\qquad \log N=o(k).
\end{equation}
In particular \(N\) and \(X\) exceed every fixed power of \(k\).

We will also use the following fixed-parameter consequence of
Shiu's short-interval estimate~\cite{Shiu}, in the form stated by
Nair and Tenenbaum~\cite[Introduction]{NT}:
\begin{equation}\label{eq:shiu}
 \sum_{j=1}^d\tau(k-j)\le C_\beta d\log k
 \quad(k^\beta\le d\le k/3).
\end{equation}
For clarity, its parameters are fixed: apply the theorem to the
multiplicative function \(\tau=\tau_2\), modulus one, and interval
\((k-d-1,k-1]\), with exponent \(\beta/2\).
This interval has length \(d\) and lies at height comparable to \(k\).
The hypotheses \(\tau(p^\nu)=\nu+1\le2^\nu\) and
\(\tau(n)\ll_\varepsilon n^\varepsilon\) hold.
The resulting bound is
\(O_\beta(d(\log k)^{-1}\exp(\sum_{p\le k}2/p))
=O_\beta(d\log k)\).
Neither the function nor the interval exponent varies with \(k\).

\begin{proposition}[Exact-center energy]\label{prop:center}
For all sufficiently large \(k\in\mathcal G\), and every reduced
fraction \(u/q\) with \(2\le q\mid k!(k+1)!\),
\begin{equation}\label{eq:center}
 \sum_{0\le n<X}\sin^2(\pi un^k/q)>k^2.
\end{equation}
Here \(X\) may be chosen as in \eqref{eq:NX} for any fixed \(B>0\).
\end{proposition}
\begin{proof}
Suppose the energy is at most \(k^2\). Begin with
\(F(x)=ux^k/q\), of degree \(d=k\).
Lemma~\ref{lem:contraction} and \(\log q=O(kK)\) give a new degree
\[
 M\le CkK^3/Z+CK^4=o(k).
\]
In particular \(M<k/3\), and \(q\mid G_M(k)\).
By \eqref{eq:newtoncriterion}, replace \(F\) by
\[
 \frac uq\sum_{j=0}^M A_j(k)\binom{x}{j}.
\]
Its difference from \(un^k/q\) is integral at every integer, so the
sine energy has not changed. The modulus \(q\) and prefix \(X\) also
remain unchanged.
Every subsequent invocation of the rounding lemma starts again on this
same prefix, so its interval-length loss is not accumulated across steps.

Whenever \(k^\beta\le d\le k/3\) and \(q\mid G_d(k)\),
\eqref{eq:logG}, \(\log D(v)\le\tau(v)\log(v+1)\), and
\eqref{eq:shiu} give
\[
 \log q\le C_\beta dK^2.
\]
A further contraction therefore gives
\[
 M\le C_\beta dK^4/Z+CK^4<d/2
\]
for large \(k\), uniformly in this range.
After finitely many steps the degree is at most \(k^\beta\).

Choose the dyadic \(Y\) for which \eqref{eq:good} holds for \(k\).
Put \(E_Y=\exp(a\mathcal L(Y))\).
For every \(1\le d\le k^\beta\), the same bound \eqref{eq:good} gives
\[
 \log q\le d\log(2k)+dE_Y.
\]
The contraction estimate now reads
\begin{equation}\label{eq:smallcontraction}
 M\le\rho_k d+C_0K^4,\qquad
 \rho_k=\frac{C_2K^2(K+E_Y)}Z\longrightarrow0.
\end{equation}
The constants are fixed. Since \(Y\le k\le2Y\),
\(\mathcal L(Y)=\mathcal L(k)+O(1/\log\log k)\).
Thus \(E_Y=\exp(a\mathcal L(k)+o(1))\), whereas
\(Z\sim\exp(c\mathcal L(k))\); the inequality \(a<c\) gives
\(\rho_k\to0\), uniformly over the selected degrees.
Take \(k\) large enough that \(\rho_k<1/4\).
While \(d>2C_0K^4\), this gives \(M<3d/4<k\), so the same exact
replacement remains valid.
Eventually \(q\mid G_d(k)\) with \(1\le d\le2C_0K^4\).
A reduction to degree zero would already give \(q\mid G_0(k)=1\),
which is impossible. At the terminal positive degree,
\begin{equation}\label{eq:smallq}
 \log q\le CK^4(K+E_Y)=o(Z).
\end{equation}
Thus \(q=X^{o(1)}\).

Return to the original phase, not to an individual cyclotomic component.
Each full block of \(q\) consecutive inputs contains an \(n\equiv1\pmod q\).
Since \(\gcd(u,q)=1\),
\[
 \sin^2(\pi u/q)\ge4\dist{u/q}^2\ge4/q^2.
\]
Consequently the original energy is at least
\[
 \frac{4\lfloor X/q\rfloor}{q^2}
 \ge\frac{2X}{q^3}=X^{1-o(1)}>k^2,
\]
by \eqref{eq:scales} and \eqref{eq:smallq}. This is the desired contradiction.
\end{proof}

\section{From rational centers to collision estimates}\label{sec:collision}

Retain \(k\in\mathcal G\) and \(N\) from \eqref{eq:NX}. Put
\[
 S=512\left\lfloor\frac{\lfloor N/k\rfloor}{512}\right\rfloor,\quad
 \mathcal V=\{N-S+1,\ldots,N\},\quad
 R=\{1,\ldots,N-S\}\setminus\{2,\ldots,257\}.
\]
Thus \(\mathcal V\) is the reserved set of bases and \(R\) is the available
initial population. Write
\[
 M_R=|R|,\qquad t=\lfloor N/8192\rfloor,\qquad \ell=128.
\]
For \(b\in\mathcal V\), let \(q_b=b^k\), and define
\begin{align*}
 D_b(\alpha)&=\sum_{n\in R}\sin^2(\pi\alpha n^k),\\
 F_b(h)&=\frac1{M_R}\sum_{n\in R}
           \cos^2(\pi hn^k/q_b)
       =1-\frac{D_b(h/q_b)}{M_R}.
\end{align*}
We will prove
\begin{equation}\label{eq:collision}
 \sup_{b\in\mathcal V}\sum_{h=1}^{q_b-1}F_b(h)^t=o(1).
\end{equation}
Notice that \(M_R\sim N\) and \(t/M_R\ge1/16384\) for large \(k\).
Also \(S\le N/k\), so \(b/N\ge1-1/k\) for \(b\in\mathcal V\).
Since \((1-1/k)^k\to\e^{-1}>1/4\), we have
\(b^k\ge N^k/4\) uniformly over the reserved bases.

\begin{lemma}[A common rational approximation]\label{lem:inverse}
Let \(A>0\) be fixed. There are a constant \(C\) and an integer
\(B_k\mid k!(k+1)!\), independent of \(\alpha\), such that
\[
 D_b(\alpha)\le Ak\log N
 \quad\Longrightarrow\quad
 \left|\alpha-\frac u{B_k}\right|
 \le N^{-k}\exp(CkK^6)
\]
for some integer \(u\). The assertion is uniform in \(b\in\mathcal V\).
\end{lemma}
\begin{proof}
The population is independent of \(b\). Add back the at most 256 excluded
terms and use the prefix \(0\le n<N-S\); the term at zero vanishes.
Its energy \(E_0\) is at most \(Ak\log N+256\).

Take \(P=P_{k+1}=(1-z)^{k+1}Q\) from Lemma~\ref{lem:flat}, and
\(M=Q(1)\mid(k+1)!\).
Round \(f(n)=\alpha n^k\) to integers \(m_n\).
As in Lemma~\ref{lem:rounding}, the number of nonzero integer residuals
\(P(E)m_n\) is at most \(\normone P^2E_0/4\).
We have \(\log(1+E_0)=O(K)\), including when \(E_0=0\).
The zero-run count in Lemma~\ref{lem:rounding} therefore provides
\(L\ge N\exp(-C_1K^6)\) consecutive valid residual positions,
say \(s,\ldots,s+L-1\).

Put \(D_0=\deg P\). As \(\deg Q=D_0-k-1\), the integers
\(b_n=Q(E)m_n\) are defined for every
\(s\le n\le s+L+k\). On the \(L\) residual positions,
\(\Delta^{k+1}b_n=0\).
Successive finite differences consequently show that all these
values agree with one rational polynomial of degree at most \(k\).
Its Newton coefficients about the integer \(s\) are integral, so
its coefficient of \(n^k\) is \(u/k!\) for some \(u\in\Z\).
The polynomial \(Q(E)f\) has coefficient \(M\alpha\) at \(n^k\).
Since \(Q\) has nonnegative coefficients of sum \(M\), the finite
convolution inequality bounds the squared error of their difference,
on the first \(L\) positions of this interval, by \(M^2E_0/4\).

Take \(H=\lfloor L/(2k)\rfloor\ge L/(4k)\).
The \(k\)th difference with step \(H\) of this error polynomial is the
constant
\[
 k!H^kM\left(\alpha-\frac u{k!M}\right).
\]
There are at least \(L/2\) valid starting points. The difference operator
has coefficient norm \(2^k\), so
\[
 \frac L2(k!H^kM)^2
 \left|\alpha-\frac u{k!M}\right|^2
 \le4^kM^2E_0/4.
\]
Using \(k!\ge(k/\e)^k\), we obtain
\[
 \left|\alpha-\frac u{k!M}\right|
 \le(8\e)^k L^{-k}\sqrt{\frac{E_0}{2L}}
 \le N^{-k}\exp(CkK^6).
\]
Take \(B_k=k!M\). All filters and this denominator depend only on \(k\).
\end{proof}

\begin{lemma}[Small symmetric numerators]\label{lem:smallnumerators}
Put \(K_* =10^{10}(1+\log N/\log k)\). Uniformly for \(b\in\mathcal V\),
\[
 \sum_{\substack{1\le h<q_b\\
 \min(h,q_b-h)\le (b/(K_*k^2))^k}}F_b(h)^t=o(1).
\]
\end{lemma}
\begin{proof}
By symmetry take \(a=\min(h,q_b-h)\) and put \(r=ba^{-1/k}\).
Then \(K_*k^2\le r\le b\le N\), and the sine phase is \((n/r)^k\).

If \(r>N/2\), use the interval
\[
 [r(1/8)^{1/k},\,r(1/4)^{1/k}].
\]
It has \(\gg N/k\) integer points, all in \(R\) for large \(k\):
its upper endpoint is at most \(N4^{-1/k}<N-S\), and its lower
endpoint exceeds 257. On it the sine square is at least \(1/8\).
Thus \(D_b(h/q_b)\gg N/k\).

If \(r\le N/2\), set \(Y_r=\lfloor r/(1000k)\rfloor\).
For \(1\le j\le Y_r\), use the disjoint intervals
\[
 I_j=[r(j+1/8)^{1/k},\,r(j+1/4)^{1/k}].
\]
The mean value theorem gives
\[
 |I_j|\ge \frac r{16k(j+1)}>2.
\]
Their integer points lie in \(R\).
Indeed their lower endpoints exceed 257, while their upper endpoints
are at most \(r(Y_r+1)^{1/k}\le3N/4\) eventually, because
\(\log N=o(k)\). Each interval contains at least half its length in
integers, with sine square at least \(1/8\). Therefore
\[
 D_b(h/q_b)
 \ge\frac r{256k}\sum_{j=1}^{Y_r}\frac1{j+1}
 \ge\frac{r\log(Y_r+1)}{512k}
 \ge\frac{K_*k\log k}{1024}.
\]
The first case satisfies the same final bound, since
\(N/k\gg K_*k\log k\).

Now \(F_b(h)^t\le\exp(-D_b(h/q_b)/16384)\).
There are fewer than \(N^k\) frequencies, and
\(K_*>2^{24}\log N/\log k\) by a fixed factor. Hence their total
contribution is bounded by
\[
 N^k\exp\!\left(-\frac{K_*k\log k}{2^{24}}\right)=o(1).
\]
\end{proof}

\begin{proposition}\label{prop:collision}
The uniform collision estimate \eqref{eq:collision} holds.
\end{proposition}
\begin{proof}
Choose a fixed \(A\) sufficiently large.
Frequencies with \(D_b(h/q_b)>Ak\log N\) have total contribution at most
\[
 N^k\exp(-Ak\log N/16384)=o(1).
\]
Every remaining frequency is in an arc
\[
 \left|\frac h{q_b}-\frac u{B_k}\right|\le\delta,\qquad
 \delta=N^{-k}\exp(CkK^6),
\]
by Lemma~\ref{lem:inverse}.
We have \(\log B_k=O(kK)\), and \(\delta B_k=o(1)\).
Thus at most \(B_k+2\) centers are relevant in or adjacent to \([0,1]\).
Each arc contains at most \(2q_b\delta+2\) points of the actual grid.
The total number of grid points in these arcs is therefore at most
\(\exp(C_3kK^6)\).
Consequently those among them with \(D_b\ge A_3kK^6\) also contribute
\(o(1)\), for a fixed \(A_3>2\cdot16384C_3\).

It remains to consider \(D_b<A_3kK^6\).
Choose the constant \(B\) in \eqref{eq:NX} larger than \(C+2\).
On the prefix \(n<X\), replacing the frequency by its center changes
the square root of the sine energy by at most
\[
 \pi\delta X^k\sqrt X
 \le \pi\exp(-(B-C)kK^6+\tfrac12\log X)=o(1).
\]
If the center is not an integer, its reduced denominator is a divisor
of \(B_k\), and Proposition~\ref{prop:center} gives center energy
greater than \(k^2\). Removing the at most 256 auxiliary terms leaves
\[
 D_b(h/q_b)\ge k^2/2-256
\]
for large \(k\), contradicting \(D_b<A_3kK^6\).

If the center is an integer, then
\(\min(h,q_b-h)\le q_b\delta\).
Since \(\log N\) dominates \(K^6\) and
\(\log(K_*k^2)=O(K)\), we have
\(\delta\le(K_*k^2)^{-k}\) eventually.
These frequencies lie in the range of Lemma~\ref{lem:smallnumerators},
whose total contribution is \(o(1)\).
All cases have now been covered.
\end{proof}

\section{A represented interval and the infinite tail}\label{sec:completion}

The last step converts \eqref{eq:collision} into ordinary subset sums,
with no repetition of bases.

\begin{lemma}\label{lem:groups}
For all sufficiently large \(k\in\mathcal G\), there are \(\ell=128\)
pairwise disjoint sets of bases, all at most \(N\), whose power
subset-sum sets \(A_1,\ldots,A_\ell\) satisfy
\begin{equation}\label{eq:groups}
 |A_i|\ge sN^k/16,\qquad A_i\subseteq[0,2sN^k],
 \qquad s=S/512.
\end{equation}
No \(A_i\) is contained in an arithmetic progression of integer
difference greater than one.
\end{lemma}
\begin{proof}
Choose \(\ell\) disjoint random \(t\)-subsets \(C_i\) of \(R\), uniformly.
This is possible since \(\ell t<M_R\), and each marginal is a uniform
\(t\)-subset. Independence between the \(C_i\) is not required.

For a fixed \(C=C_i\), choose each subset \(U\subseteq C\) with
probability \(2^{-t}\); equivalently, include its bases independently
with probability \(1/2\). Let
\[
 \mu_b(x)=2^{-t}\#\left\{U\subseteq C:
             \sum_{n\in U}n^k\equiv x\pmod {q_b}\right\}.
\]
This is a distribution on residues, not the uniform distribution
on distinct represented values. Its Fourier transform is
\[
 \widehat\mu_b(h)=
 \prod_{n\in C}\frac{1+\exp(2\pi\ii hn^k/q_b)}2.
\]
Finite Fourier orthogonality therefore gives the collision excess
\[
 q_b\sum_x\mu_b(x)^2-1
 =\sum_{h=1}^{q_b-1}\prod_{n\in C}\cos^2(\pi hn^k/q_b).
\]
For \(x_n=\cos^2(\pi hn^k/q_b)\), the expectation of the product over
a uniform \(t\)-subset equals \(e_t((x_n)_{n\in R})/\binom{M_R}{t}\),
where \(e_t\) is the elementary symmetric polynomial of degree \(t\).
Maclaurin's inequality bounds this by
\((M_R^{-1}\sum_{n\in R}x_n)^t=F_b(h)^t\).

Write \(\eta_k=\sup_{b\in\mathcal V}\sum_{h=1}^{q_b-1}F_b(h)^t=o(1)\).
If the support has fewer than \(q_b/4\) elements, Cauchy--Schwarz gives
\(q_b\sum_x\mu_b(x)^2>4\), so the collision excess is greater than three.
For each fixed group, the expected number of such bad reserved bases
is at most \(S\eta_k/3\). Markov's inequality bounds the probability
of more than \(S/2\) bad bases by \(2\eta_k/3\).
A union bound over the fixed \(\ell\) groups proves that, with
probability \(1-o(1)\), every group has at least \(S/2\) good bases.

We may simultaneously require
\begin{equation}\label{eq:initialspan}
 \sum_{n\in C_i}n^k\le4tN^k/k\qquad(1\le i\le\ell).
\end{equation}
Indeed the expectation is at most \(2tN^k/k\), because
\(\sum_{n\le N}n^k\le N^{k+1}/(k+1)+N^k\) and \(M_R\sim N\).
Sampling without replacement has variance at most \(tN^{2k}\).
Chebyshev's inequality makes the failure probability at most
\(k^2/(4t)=o(1)\), uniformly in the fixed number of groups.
Fix a realization with both properties.

Allocate \(s\) good reserved bases to each group, with no base used
twice. Before any individual choice, fewer than
\(\ell s=S/4\) bases have been used, whereas the group in question
has at least \(S/2\) good choices. The greedy allocation is therefore
possible; the completed allocation uses exactly \(S/4\) bases.
If a finite integer set \(A\) occupies \(r\) residue classes modulo
a positive integer \(q\), then
\[
 |A+\{0,q\}|\ge|A|+r.
\]
For each occupied residue, its largest element supplies a new element
after addition of \(q\).
Applying this with \(q=b^k\) shows that each allocated base increases
the subset-sum cardinality by at least \(q_b/4\ge N^k/16\).
The original support is retained after earlier additions, so the
support condition continues to hold. After \(s\) additions the
cardinality is at least \(sN^k/16\).

Finally add the distinct auxiliary pairs
\[
 (2,3),(4,5),\ldots,(256,257)
\]
to the 128 groups, one pair per group.
They are disjoint from \(R\) and from the reserved bases.
Each resulting subset-sum set contains zero and two coprime \(k\)th
powers, so it cannot lie in a proper integer arithmetic progression.
The initial span in \eqref{eq:initialspan} is at most \(sN^k/3\)
eventually, because \(4t/(ks)\to1/4<1/3\).
The reserved bases add at most \(sN^k\), and the two auxiliary powers
add at most \(2\cdot257^k=o(sN^k)\).
This proves \eqref{eq:groups}.
\end{proof}

We use Lev's consecutive-interval theorem in the following sufficient
form~\cite[Corollary~1]{Lev}. If \(A_1,\ldots,A_\ell\) are subsets of
\([0,l]\cap\Z\), each of size at least \(n\ge3\), none contained in an
integer arithmetic progression of difference greater than one, and
\[
 \ell\ge2\left\lceil\frac{l-1}{n-2}\right\rceil,
\]
then their sumset contains \(\ell(n-1)+1\) consecutive integers.
Here the length of a block means its largest element minus its smallest,
as in Lev's statement.

\begin{proof}[Proof of Proposition~\ref{prop:upper}]
Apply this result to Lemma~\ref{lem:groups}, with
\[
 l=\lceil2sN^k\rceil,\qquad n=\lfloor sN^k/16\rfloor,\qquad \ell=128.
\]
The ratio \((l-1)/(n-2)\) tends to 32, so the strict inequality has
a fixed margin. The sumset contains an interval \([a,a+L]\cap\Z\)
with
\[
 L\ge128(n-1)\sim8sN^k>(N+1)^k.
\]
Also \(0\le a\le256sN^k\le256N^{k+1}\).
Every integer in this interval is a sum of distinct \(k\)th powers
with bases at most \(N\), because the groups of bases are disjoint.

Adjoin the unused powers with bases \(N+1,N+2,\ldots\).
If the currently represented interval is \([a,a+L]\) and a new power
\(b^k\) satisfies \(b^k\le L+1\), its translate joins it without a gap,
producing \([a,a+L+b^k]\).
Initially \(L\ge(N+1)^k\).
If \(L\ge b^k\), then the new length is at least \(2b^k>(b+1)^k\),
since \(b\ge N\) and \(N/k\to\infty\).
This preserves the induction invariant.
At each stage the interval consists of subset sums of the finitely
many bases used up to that stage, all distinct. Its length tends to
infinity. Every fixed integer at least \(a\) consequently belongs
to one finite-stage interval, and is represented by finitely many
distinct positive \(k\)th powers. We have proved
\[
 T_k\le256N^{k+1}.
\]
Taking logarithms and using \eqref{eq:scales} gives \eqref{eq:upper}.
\end{proof}

\begin{proof}[Proof of Theorem~\ref{thm:main}]
If there were only finitely many strict adjacent decreases, then
\(T_{k+1}\ge T_k\) for all sufficiently large \(k\).
Choose arbitrarily large \(j\) from Lemma~\ref{lem:lower}.
Let \(Y\) be a dyadic integer with \(j\le Y<2j\), and choose
\(k\in\mathcal G\cap[Y,2Y]\).
Then \(j\le k<4j\), and
\(\mathcal L(k)=\mathcal L(j)+o(1)\).
Proposition~\ref{prop:upper} gives
\[
 \log T_k\le(4+o(1))j\exp\!\left(\frac25\mathcal L(j)\right),
\]
whereas
\[
 \log T_j\ge j\exp\!\left(\frac9{20}\mathcal L(j)\right).
\]
The latter is larger for all sufficiently large \(j\), contradicting
\(T_j\le T_k\). Thus strict adjacent decreases occur infinitely often.
\end{proof}

\section*{Acknowledgment}
The authors acknowledge extensive assistance from OpenAI's
artificial-intelligence tools in developing and checking the
mathematical argument and in preparing the exposition.

\begin{thebibliography}{9}
\bibitem{BK}
A.~S.~Belov and S.~V.~Konyagin,
\emph{An estimate of the free term of a non-negative trigonometric
polynomial with integer coefficients},
Izvestiya: Mathematics \textbf{60} (1996), no.~6, 1123--1182.
\href{https://doi.org/10.1070/IM1996v060n06ABEH000095}
{doi:10.1070/IM1996v060n06ABEH000095}.

\bibitem{EG}
P.~Erd\H{o}s and R.~L.~Graham,
\emph{Old and New Problems and Results in Combinatorial Number Theory},
Monographies de L'Enseignement Math\'ematique, vol.~28,
Universit\'e de Gen\`eve, Geneva, 1980.
\href{https://fanchung.ucsd.edu/ron/papers/80_11_number_theory.pdf}{Author's archived copy}.

\bibitem{FP}
Kai (Steve) Fan and Paul Pollack,
\emph{The maximal order of the shifted-prime divisor function},
Integers \textbf{26A} (2026),
\href{https://math.colgate.edu/~integers/aap9/aap9.pdf}{Article A9}.
\href{https://doi.org/10.5281/zenodo.23017984}{doi:10.5281/zenodo.23017984}.

\bibitem{Kim}
D.~Kim,
\emph{On the largest integer that is not a sum of distinct positive
\(n\)th powers},
Journal of Integer Sequences \textbf{20} (2017), Article 17.7.5.
\href{https://cs.uwaterloo.ca/journals/JIS/VOL20/Kim/kim6.pdf}{Journal version}.

\bibitem{Lev}
V.~F.~Lev,
\emph{Consecutive integers in high-multiplicity sumsets},
Acta Mathematica Hungarica \textbf{129} (2010), 245--253.
\href{https://doi.org/10.1007/s10474-010-0026-6}{doi:10.1007/s10474-010-0026-6}.
\href{https://arxiv.org/abs/0806.4580v1}{arXiv:0806.4580v1}.

\bibitem{NT}
M.~Nair and G.~Tenenbaum,
\emph{Short sums of certain arithmetic functions},
Acta Mathematica \textbf{180} (1998), 119--144.
\href{https://tenenb.perso.math.cnrs.fr/PPP/ShortSums.pdf}{Author's version}.

\bibitem{Shiu}
P.~Shiu,
\emph{A Brun--Titchmarsh theorem for multiplicative functions},
Journal f\"ur die reine und angewandte Mathematik \textbf{313} (1980), 161--170.
\href{https://doi.org/10.1515/crll.1980.313.161}{doi:10.1515/crll.1980.313.161}.

\end{thebibliography}
\end{document}
